\hnheader[Erst diagnostizieren, dann reparieren -- nicht planlos herumprobieren.][Lernkurve lesen: Lücke $\Rightarrow$ Varianz, hoher Train-Fehler $\Rightarrow$ Bias]{Praktische}{Fehleranalyse}{Diagnostik, Error \& Ablative Analysis}
\hnsrc{materials/ML-advice.pdf (A.\ Ng), notes/error-analysis.pdf (Y.\ Le Calonnec), section/error-analysis.pdf}

\begin{hngrid}
%
\begin{hnp}[raster multicolumn=2,acc=hnorange]{1}{Problem}
Spamfilter mit $20\%$ Testfehler. Mögliche Maßnahmen: mehr Daten, weniger/mehr Features, mehr Iterationen, Newton, anderes $\lambda$, SVM\dots\ Blindes Ausprobieren kostet Wochen. \hlt{Diagnose zuerst!}
\end{hnp}
%
\begin{hnp}[raster multicolumn=2,acc=hnpurple]{2}{Lernkurven}
Plotte Train- und Testfehler über der Trainingsgröße $m$:
\begin{itemize}
\item \textbf{Varianz}: große Lücke, Test fällt noch $\Rightarrow$ mehr Daten helfen
\item \textbf{Bias}: Train schon hoch, kleine Lücke $\Rightarrow$ Modell ausdrucksstärker machen
\end{itemize}
\end{hnp}
%
\begin{hnp}[raster multicolumn=2,acc=hnteal]{3}{Skizze: zwei Diagnosen}
\centering
\begin{tikzpicture}[>=Stealth,scale=.62]
  \begin{scope}
    \draw[->,hnnavy] (0,0)--(3.6,0); \draw[->,hnnavy] (0,0)--(0,3);
    \draw[hnorange,thick] (.2,2.6) .. controls (1,1.9) and (2,1.6) .. (3.4,1.4);
    \draw[hnpurple,thick] (.2,.3) .. controls (1,.5) and (2,.6) .. (3.4,.6);
    \draw[dashed,gray] (0,1.0)--(3.6,1.0);
    \node[font=\tiny] at (1.8,-.35) {hohe Varianz};
    \node[font=\tiny,hnorange] at (3,1.85) {Test};
    \node[font=\tiny,hnpurple] at (3,.25) {Train};
  \end{scope}
  \begin{scope}[shift={(4.6,0)}]
    \draw[->,hnnavy] (0,0)--(3.6,0); \draw[->,hnnavy] (0,0)--(0,3);
    \draw[hnorange,thick] (.2,2.8) .. controls (1,2.3) and (2,2.1) .. (3.4,2.0);
    \draw[hnpurple,thick] (.2,1.2) .. controls (1,1.6) and (2,1.8) .. (3.4,1.85);
    \draw[dashed,gray] (0,1.0)--(3.6,1.0);
    \node[font=\tiny] at (1.8,-.35) {hoher Bias};
  \end{scope}
\end{tikzpicture}\\[-2pt]
{\tiny gestrichelt: gewünschtes Niveau}
\end{hnp}
%
\begin{hnp}[raster multicolumn=3,acc=hnnavy]{4}{Welche Maßnahme behebt was?}
\renewcommand{\arraystretch}{1.22}
\begin{tabular}{@{}p{4.0cm}p{4.4cm}@{}}
\hline
\textbf{Maßnahme}&\textbf{behebt}\\\hline
mehr Trainingsbeispiele&hohe Varianz\\
weniger Features&hohe Varianz\\
mehr / bessere Features&hohen Bias\\
mehr Iterationen, Newton&Optimierungs\emph{algorithmus}\\
anderes $\lambda$ / $C$, andere Klasse&Optimierungs\emph{ziel}\\\hline
\end{tabular}
\end{hnp}
%
\begin{hnp}[raster multicolumn=3,acc=hngreen]{5}{Algorithmus oder Ziel?}
Ein SVM schlägt BLR (Bayessche Logistische Regression) bei der wahren Metrik $a(\theta)$, obwohl BLR das Ziel $J(\theta)$ maximiert ($\theta$ Parameter, $J$ Zielfunktion, $a$ Metrik, die uns wirklich interessiert). Vergleiche $J(\theta_{\mathrm{SVM}})$ mit $J(\theta_{\mathrm{BLR}})$:
\begin{pcode}[Diagnose]
\kw{if} $J(\theta_{\mathrm{BLR}})<J(\theta_{\mathrm{SVM}})$:\\
\hii Problem = \kw{Optimierer} (nicht konvergiert)\\
\hii $\to$ mehr Iterationen, Newton, andere $\alpha$\\
\kw{else}: \ \cm{\# $J$ erreicht, aber $a$ schlecht}\\
\hii Problem = \kw{Zielfunktion} falsch\\
\hii $\to$ anderes $\lambda/C$, Gewichte, andere Klasse
\end{pcode}
\end{hnp}
%
\begin{hnp}[raster multicolumn=3,acc=hnrose]{6}{Error Analysis (Pipeline)}
Bei mehrstufigen Systemen jede Komponente einzeln durch die \emph{Ground Truth} ersetzen und die Gesamtgenauigkeit messen:
\renewcommand{\arraystretch}{1.15}
\begin{tabular}{@{}lc@{}}
\hline
+ perfektes Preprocessing&85.1\,\%\\
+ perfekte Gesichtsdetektion&91\,\%\\
+ perfekte Augen-Segmentierung&95\,\%\\\hline
\end{tabular}\\
(Start: 85\,\%.) Große Sprünge zeigen, wo Arbeit sich lohnt.
\end{hnp}
%
\begin{hnp}[raster multicolumn=3,acc=hnpurple]{7}{Ablative Analyse}
Umgekehrt: vom \emph{guten} System aus Komponenten entfernen.
\renewcommand{\arraystretch}{1.15}
\begin{tabular}{@{}lc@{}}
\hline
Gesamtsystem&99.9\,\%\\
$-$ Rechtschreibung&99.0\,\%\\
$-$ Text-Parser&95\,\%\\
$-$ Bild-Features (Baseline)&94\,\%\\\hline
\end{tabular}\\
Zeigt, welche Features wirklich beitragen (gut für Publikationen).
\end{hnp}
%
\begin{hnex}[raster multicolumn=3]{Beispiel: Fehler manuell zählen}
Katzenklassifikator, Accuracy $90\%$. Von 100 Fehlern sind \ldots\\
\hnsol\ \textbf{5 Hunde}: Selbst mit perfektem Hunde-Filter steigt die Accuracy nur $90\%\to\ora{90.5\%}$ -- nicht lohnend.\\
\textbf{50 Hunde}: $\to\ora{95\%}$ -- \hlm{hier lohnt sich die Arbeit!}
\end{hnex}
%
\begin{hnp}[raster multicolumn=3,acc=hnteal]{8}{Vorgehen in 5 Schritten}
\begin{enumerate}
\item Baseline bauen, Metrik festlegen
\item Lernkurve $\to$ Bias oder Varianz?
\item Optimierer vs.\ Ziel prüfen
\item Error Analysis der Pipeline
\item ca.\ 100 Fehler kategorisieren, priorisieren
\end{enumerate}
\end{hnp}
%
\begin{hnp}[raster multicolumn=6,acc=hnorange]{9}{Key Takeaways}
\begin{hnchk}
\item Diagnose vor Aktion
\item Lernkurve: Lücke = Varianz, Train hoch = Bias
\item $J$ vs.\ $a$: Optimierer oder Ziel schuld?
\item Fehleranalyse zeigt die Obergrenze des Gewinns
\end{hnchk}
\end{hnp}
%
\begin{hnp}[raster multicolumn=6,acc=hnpurple,colback=hnlilac!30]{?}{Selbsttest: kannst du das erklären?}
\hnquiz{Lernkurve: Varianz oder Bias?}{Große Lücke $\Rightarrow$ Varianz; hoher Trainingsfehler $\Rightarrow$ Bias.}{$J(\theta_{\mathrm{BLR}})<J(\theta_{\mathrm{SVM}})$ bedeutet?}{Optimierer hat sein Ziel nicht erreicht (Konvergenzproblem).}{Error vs.\ Ablative Analysis?}{Error: Komponenten durch Ground Truth ersetzen; Ablation: Komponenten entfernen.}
\end{hnp}
%
\begin{hnbanner}[raster multicolumn=6]
„Wer nicht misst, verbessert das Falsche.“
\end{hnbanner}
\end{hngrid}
